\chapter{PHÂN TÍCH MÃ NGUỒN VÀ CHỨC NĂNG}

Mã nguồn được tổ chức thành một gói Python \texttt{churn} và các tệp kiểm thử; chi tiết cài đặt và chạy nằm trong \texttt{README.md}.

\section{Cấu trúc các module}
\begin{longtable}{p{4.2cm} p{10cm}}
    \caption{Các module chính của dự án} \label{tab:modules} \\
    \hline \textbf{Tệp} & \textbf{Chức năng} \\ \hline
    \endfirsthead
    \hline \textbf{Tệp} & \textbf{Chức năng (tiếp)} \\ \hline
    \endhead
    \texttt{config.py} & Tham số tái lập (\texttt{RANDOM\_STATE}), đường dẫn và danh sách cột. \\
    \texttt{data.py} & Nạp CSV, làm sạch, thống kê chất lượng dữ liệu. Không học tham số. \\
    \texttt{features.py} & Đặc trưng kỹ thuật (Luồng B) và bộ tiền xử lý trong \texttt{Pipeline}. \\
    \texttt{models.py} & Định nghĩa baseline, K-NN, Naive Bayes, Cây quyết định và lưới tham số. \\
    \texttt{evaluate.py} & Chia train/test một lần, \texttt{GridSearchCV} trên tập train, tính các độ đo trên tập test. \\
    \texttt{error\_analysis.py} & Tỷ lệ sai theo nhóm và danh sách các ca bỏ sót (FN). \\
    \texttt{run\_all.py} & Chạy toàn bộ thí nghiệm, xuất bảng (CSV, LaTeX), hình và \texttt{run\_info.json}. \\
    \texttt{tests/test\_pipeline.py} & Kiểm thử làm sạch, chống rò rỉ và tái lập. \\
\end{longtable}

\section{Luồng xử lý khi chạy}
Lệnh \texttt{python -m churn.run\_all} thực hiện tuần tự: nạp dữ liệu, làm sạch, vẽ hình EDA, chia train/test,
chạy thí nghiệm cho hai luồng đặc trưng và ba thuật toán cùng baseline, chọn cấu hình tốt nhất theo CV,
vẽ hình đánh giá, phân tích lỗi và ghi kết quả vào thư mục \texttt{results/}.

\section{Đoạn mã quan trọng: chống rò rỉ dữ liệu}
Đoạn mã dưới đây thể hiện nguyên tắc cốt lõi: toàn bộ bước học tham số nằm trong Pipeline và được fit trên tập train.
\begin{verbatim}
# Pipeline: bộ tiền xử lý + mô hình, cùng được fit trên tập train
pipe = make_pipeline(model, X_tr)
search = GridSearchCV(pipe, grid, cv=cv, scoring="f1")
search.fit(X_tr, y_tr)        # chỉ dữ liệu huấn luyện
# Tập kiểm thử chỉ dùng để đánh giá, không dùng để chọn
y_pred = search.best_estimator_.predict(X_te)
\end{verbatim}
